\section{Basic Attention：从固定压缩到按需读取}

在进入 self-attention 前，先掌握编码器—解码器 attention。编码器为每个源位置产生一组状态 $h_1,\ldots,h_n$。在时间 $t$，decoder 使用状态 $s_{t-1}$ 计算兼容度、归一化权重和 context：
\[
e_{t,i}=v^\top\tanh(W_hh_i+W_ss_{t-1}),\qquad
\alpha_{t,i}=\frac{\exp(e_{t,i})}{\sum_j\exp(e_{t,j})},\qquad
c_t=\sum_i\alpha_{t,i}h_i.
\]
其中 $e_{t,i}$ 是第 $t$ 个 query 与源位置 $i$ 的 score，$\alpha_{t,i}$ 是归一化 attention weight，$c_t$ 是本步读取结果。这个 Bahdanau-style additive attention 把一个 context vector 改成每步不同的 context。

\subsection{Soft 与 hard attention}

\term{soft attention} 对所有候选位置使用连续权重，通常可通过反向传播端到端训练；\term{hard attention} 选择或采样少数位置，计算可能更稀疏，但离散选择通常不可直接微分，需要 policy-gradient、relaxation 或其他 estimator。二者不是“软更准确、硬更快”的固定排序，真实效果取决于任务和训练方法。

\lecturefigure{slide-09-part-soft-attention.jpg}{早期视觉 attention 已区分 soft 与 hard selection。}{官方视频 06:30；Xu et al. 2015 slide citation。}

\begin{knowledgebox}{读图：连续权重与离散选择改变训练接口}
Soft attention 像对多个区域做可微加权平均，梯度能流向每个位置；hard attention 像选择一个或少数区域，计算路径更明确，却引入采样方差和信用分配问题。图中的 cat 示例不是为了证明模型“看懂猫”，而是展示注意力权重可作为信息路由的中间变量。权重可视化也不自动等于因果解释。
\end{knowledgebox}

\begin{warningbox}{Attention weight 不必等于人类解释}
高权重表示当前计算读取了某个 value 的程度，但模型还经过 residual、multi-head、nonlinearity 和后续层。不能仅凭一张 heatmap 断言模型的完整推理原因；需要 ablation、counterfactual 或 causal intervention 验证。
\end{warningbox}

\subsection{Global 与 local attention}

\term{global attention} 在每个 decoding step 对全部源位置评分；\term{local attention} 只在预测窗口或选定邻域内评分。Global 提供完整访问但成本随长度增长，local 引入近邻偏置并减少计算，却可能错过远距离证据。后来的 sliding-window、block-sparse 与 routing attention 都可以看成这类设计空间的延伸。

\lecturefigure{slide-10-local-global-attention.jpg}{Global attention 读取全序列，local attention 只读取局部窗口。}{官方视频 07:10；Luong et al. 2015。}

\begin{knowledgebox}{读图：窗口是归纳偏置，不只是优化技巧}
Global model 假设任意位置都可能相关；local model 假设大多数依赖集中在邻域。窗口大小、中心预测方式与边界处理会改变可表达关系。对于长文档、音频或视频，局部结构往往真实存在，但跨段引用也同样重要；系统需要在吞吐、内存和任务依赖之间选择，而不是把 sparse attention 当无损替代。
\end{knowledgebox}

\begin{importantbox}{Attention 的第一次接口重构}
早期 attention 仍保留 recurrent encoder/decoder，但把读取方式从 fixed vector 改为 content-dependent weighted retrieval。Self-attention 的下一步，是让同一序列中的每个位置都通过类似接口读取其他位置。
\end{importantbox}

\subsection{本章小结}

Basic attention 已经包含后来系统的核心设计轴：连续或离散、全局或局部、完整访问或稀疏路由。Transformer 的创新不是发明“权重”，而是把 attention 提升为主要序列计算层。

